% ================================================================
% PREAMBLE - CẤU HÌNH LATEX CHO TUYỂN TẬP ĐỀ THI
% Đại học Sài Gòn - Khoa Toán Ứng dụng
% ================================================================

% ----------------------------------------------------------------
% DOCUMENT CLASS VÀ ENCODING
% ----------------------------------------------------------------
\documentclass[12pt, a4paper]{article}
\usepackage[utf8]{inputenc}
\usepackage[T5]{fontenc}
\usepackage{vntex}                    % Hỗ trợ tiếng Việt tốt hơn

% ----------------------------------------------------------------
% PACKAGES TOÁN HỌC
% ----------------------------------------------------------------
\usepackage{amsmath, amssymb, amsthm}
\usepackage{mathtools}                % Mở rộng amsmath
\usepackage{bm}                       % Bold math symbols

% ----------------------------------------------------------------
% LAYOUT VÀ GEOMETRY
% ----------------------------------------------------------------
\usepackage{geometry}
\geometry{
    a4paper,
    left=2.5cm,
    right=2.5cm,
    top=3cm,
    bottom=3cm,
    headheight=15pt
}

% ----------------------------------------------------------------
% PACKAGES FORMATTING
% ----------------------------------------------------------------
\usepackage{enumitem}                 % Tùy chỉnh danh sách
\usepackage{fancyhdr}                 % Header và footer
\usepackage{xcolor}                   % Màu sắc
\usepackage{framed}                   % Khung cho câu hỏi
\usepackage{titlesec}                 % Tùy chỉnh tiêu đề
\usepackage{tocloft}                  % Tùy chỉnh mục lục
\usepackage[skins,breakable]{tcolorbox} % Khung màu chuyên nghiệp
\definecolor{darkred}{RGB}{153, 0, 0} % Màu đỏ sẫm chuyên nghiệp

% ----------------------------------------------------------------
% UNICODE CHARACTERS
% ----------------------------------------------------------------
\usepackage{newunicodechar}
\newunicodechar{✓}{\ensuremath{\checkmark}}
\newunicodechar{●}{\ensuremath{\bullet}}
\newunicodechar{→}{\ensuremath{\rightarrow}}
\newunicodechar{⇒}{\ensuremath{\Rightarrow}}
\newunicodechar{≤}{\ensuremath{\leq}}
\newunicodechar{≥}{\ensuremath{\geq}}

% ----------------------------------------------------------------
% HEADER VÀ FOOTER
% ----------------------------------------------------------------
\pagestyle{fancy}
\fancyhf{}
\cfoot{\thepage}
\renewcommand{\headrulewidth}{0pt}
\renewcommand{\footrulewidth}{0pt}

% ----------------------------------------------------------------
% ĐỊNH NGHĨA MÀU SẮC
% ----------------------------------------------------------------
\definecolor{darkblue}{RGB}{0, 51, 102}
\definecolor{darkgreen}{RGB}{0, 102, 51}
\definecolor{darkred}{RGB}{153, 0, 0}
\definecolor{lightgray}{RGB}{245, 245, 245}
\definecolor{mediumgray}{RGB}{128, 128, 128}

% ----------------------------------------------------------------
% THEOREM STYLES
% ----------------------------------------------------------------
\newtheoremstyle{solutionstyle}
    {6pt}                             % Space above
    {6pt}                             % Space below
    {}                                % Body font
    {}                                % Indent amount
    {\bfseries}                       % Theorem head font
    {:}                               % Punctuation after theorem head
    {0.5em}                           % Space after theorem head
    {}                                % Theorem head spec

\theoremstyle{solutionstyle}
\newtheorem*{solution}{Lời giải}

% ----------------------------------------------------------------
% CUSTOM COMMANDS
% ----------------------------------------------------------------
% Phân số display-style
\newcommand{\dfrag}[2]{\dfrac{#1}{#2}}

% Khung cho câu hỏi
\newenvironment{questionbox}{%
    \begin{framed}
    \noindent\textbf{Đề bài:}\\[0.2cm]
}{%
    \end{framed}
}

% Khung cho ghi chú
\newenvironment{notebox}{%
    \begin{framed}
    \noindent\textbf{Ghi chú:}\\[0.2cm]
}{%
    \end{framed}
}

% ----------------------------------------------------------------
% CỠ CHỮ ĐỒNG BỘ 13PT
% ----------------------------------------------------------------
\usepackage{anyfontsize}
\renewcommand{\normalsize}{\fontsize{13pt}{16pt}\selectfont}
\renewcommand{\small}{\fontsize{13pt}{16pt}\selectfont}
\renewcommand{\footnotesize}{\fontsize{13pt}{16pt}\selectfont}
\renewcommand{\large}{\fontsize{13pt}{16pt}\selectfont}
\renewcommand{\Large}{\fontsize{13pt}{16pt}\selectfont}
\renewcommand{\LARGE}{\fontsize{13pt}{16pt}\selectfont}
\renewcommand{\huge}{\fontsize{13pt}{16pt}\selectfont}
\renewcommand{\Huge}{\fontsize{13pt}{16pt}\selectfont}

% ----------------------------------------------------------------
% SECTION FORMATTING
% ----------------------------------------------------------------
\titleformat{\section}
    {\bfseries}
    {\thesection}
    {1em}
    {}
    [\vspace{0.3cm}]

\titleformat{\subsection}
    {\bfseries}
    {\thesubsection}
    {1em}
    {}
    [\vspace{0.2cm}]

\titleformat{\subsubsection}
    {\bfseries}
    {\thesubsubsection}
    {1em}
    {}

% ----------------------------------------------------------------
% TABLE OF CONTENTS FORMATTING
% ----------------------------------------------------------------
\renewcommand{\cfttoctitlefont}{\bfseries}
\renewcommand{\cftpartfont}{\bfseries}
\renewcommand{\cftsecfont}{\bfseries}

% ----------------------------------------------------------------
% SPACING VÀ INDENTATION
% ----------------------------------------------------------------
\setlength{\parindent}{0pt}           % Không thụt đầu dòng
\setlength{\parskip}{6pt}             % Khoảng cách giữa các đoạn
\linespread{1.2}                      % Line spacing

% ----------------------------------------------------------------
% ENUMERATION SETTINGS
% ----------------------------------------------------------------
\setlist[enumerate]{
    leftmargin=*,
    itemsep=3pt,
    parsep=0pt,
    topsep=6pt
}

\setlist[itemize]{
    leftmargin=*,
    itemsep=3pt,
    parsep=0pt,
    topsep=6pt
}

% ----------------------------------------------------------------
% MATH SETTINGS
% ----------------------------------------------------------------
\allowdisplaybreaks                   % Cho phép ngắt trang trong align
\numberwithin{equation}{section}      % Đánh số phương trình theo section

% ----------------------------------------------------------------
% HYPERREF (nếu cần)
% ----------------------------------------------------------------
% \usepackage{hyperref}
% \hypersetup{
%     colorlinks=true,
%     linkcolor=darkblue,
%     urlcolor=darkblue,
%     citecolor=darkgreen
% }
